\chapter{Introduction}
\label{ch:introduction}

\section{Motivation}
\label{sec:intro_motivation}

Sequential decision-making under uncertainty is a defining problem of
machine intelligence: an agent must act now, on incomplete information,
to influence outcomes it will only observe later. Reinforcement learning
(RL) is the branch of machine learning built for exactly this setting.
Its successes are among the best known in the field --- backgammon played
at expert level from self-play \cite{tesauro1995}, Atari games mastered
from raw pixels \cite{mnih2015}, and the game of Go, long considered out
of reach, won against a world champion \cite{silver2016}. In each case
the recipe is the same: an agent tries actions, observes consequences,
receives a numerical score, and gradually shifts its behaviour toward
whatever earns more of that score.

In everyday terms, reinforcement learning is how one might train a pet:
nobody hands the animal a manual, but rewarded behaviour is repeated and
unrewarded behaviour fades. The technical content of this dissertation
is, at heart, a precise answer to two questions that this analogy leaves
open: \emph{what exactly should the learner be shown} (the state),
and \emph{what exactly should count as the reward} --- because, as our
results will demonstrate, the answers determine everything the agent
ultimately does.

Financial trading is one of RL's most natural --- and most unforgiving
--- testbeds. A trading agent observes market and portfolio information,
chooses among actions with immediate costs and delayed consequences, and
receives an unambiguous scalar verdict: the change in its wealth. The
mapping onto the RL formalism is almost mechanical, which is why a large
applied literature exists \cite{moody2001, deng2017, liu2020finrl}.

Trading is also unforgiving for a subtler reason, and this reason shapes
the entire methodology of the present work. In games, the correct outcome
is known: a CartPole balancer either balances or does not; a Go program
either wins or loses. In markets, no oracle labels the right action, so a
\emph{broken learning implementation can masquerade as a mediocre trading
strategy indefinitely} --- there is no error message, only a plausible-looking
equity curve. Worse, decades of financial economics
\cite{fama1970} warn that beating a simple passive strategy using public
price history alone should be rare; an experiment that appears to do so
easily is more likely to contain a bug or an unfair comparison than a
discovery \cite{bailey2014}. This dissertation therefore adopts an
\emph{understanding-first} methodology: derive each learning rule from
first principles, implement it from scratch, validate it on games with
known outcomes, and only then apply it to market data --- with baselines
designed so that honest conclusions are detectable.

\section{Problem statement}
\label{sec:intro_problem}

We study single-asset trading formulated as an episodic Markov Decision
Process (MDP): one calendar month of daily prices is one episode; the
agent begins each episode holding only cash and must end it in cash. At
each trading day it observes a \emph{complete} state --- nine numbers
summarising the market's recent behaviour, the time remaining in the
episode, and its own balance sheet --- and chooses one of three actions:
hold, buy a fixed slice of the asset, or sell a slice. Its reward at each
step is the change in its total wealth, with realistic transaction fees
deducted. The learning problem is to find, from experience alone, the
decision rule that maximises expected end-of-month wealth.

Two canonical algorithm families are implemented from scratch and
compared under identical conditions: the Monte-Carlo policy gradient
(REINFORCE~\cite{williams1992}), which learns the decision rule directly,
and value-based Deep Q-learning \cite{mnih2015}, which learns the value
of each action and acts greedily on those values. A third row --- PPO
\cite{schulman2017ppo}, the stabilised modern policy gradient, via a
reference library implementation \cite{raffin2021sb3} --- completes the
algorithm ablation.

\section{Objectives}
\label{sec:intro_objectives}

\begin{enumerate}
  \item[\textbf{O1}] Formulate single-asset trading as an episodic MDP
        with a complete observable state, a divisible-asset action model,
        and a wealth-change reward with explicit transaction costs.
  \item[\textbf{O2}] Derive, from first principles, the policy-gradient
        (REINFORCE) learning rule --- including why the implemented loss
        contains a log-probability term --- and the Deep Q-learning rule,
        and implement both from scratch in PyTorch.
  \item[\textbf{O3}] Validate both implementations on control tasks with
        known outcomes (CartPole, Flappy Bird, LunarLander) before any
        market experiment.
  \item[\textbf{O4}] Evaluate both methods on held-out real market data
        against properly capitalised passive and random baselines, under
        a chronological split, matched budgets and multiple seeds.
  \item[\textbf{O5}] Analyse the failure and success modes observed ---
        in particular reward-induced policy collapse to buy-and-hold ---
        and articulate the limitations of the formulation, including
        partial observability and non-IID data.
\end{enumerate}

These objectives were agreed through supervision and deliberately weight
\emph{depth of understanding} over breadth of features: a smaller system
in which every equation, every line of code and every reported number can
be defended, rather than a larger system assembled from library calls
whose behaviour cannot be explained.

\section{Contributions}
\label{sec:intro_contributions}

\begin{itemize}
  \item A fully derived, from-scratch implementation of REINFORCE and
        Deep Q-learning for a masked, divisible-asset trading MDP, with
        every loss line traceable to a stated equation (the traceability
        map is Appendix~\ref{app:trace}).
  \item A complete 9-feature state design with explicit rationale, and a
        pilot study demonstrating empirically why minimal states induce
        degenerate policies.
  \item An honest evaluation framework --- fully invested buy-and-hold
        baseline, chronological splits, fee-inclusive rewards, accounting
        identity verified per episode --- under which the central finding
        is established: under a risk-neutral wealth reward, the policy
        gradient converges to buy-and-hold exactly, while Deep Q-learning
        finds a nearby policy with measurably better risk-adjusted
        characteristics (lower drawdown, higher Sharpe ratio) at
        negligible return cost.
  \item A three-algorithm ablation (REINFORCE, Deep Q, PPO) and a fee
        sensitivity study showing that the value-based agent adapts its
        trading intensity to transaction costs, while the clipped policy
        gradient is bimodal across seeds --- both behaviours explained
        mechanistically from the update rules.
\end{itemize}

\section{Scope and non-goals}
\label{sec:intro_scope}

It is as important to state what this project is \emph{not}. It is not an
attempt to produce a profitable trading system: the efficient-market
null hypothesis \cite{fama1970}, the single asset, and the daily-bar
granularity all argue against that reading, and Chapter~\ref{ch:results}
shows why apparent profitability claims in this setting demand suspicion
before celebration. It is also not a survey of every RL algorithm:
methods are included only where they earn their place in the ablation
logic (REINFORCE as the transparent foundation, Deep Q as the
structurally different value-based family, PPO as the stabilised
descendant). The deliverable is a defensible scientific account: a
carefully formulated problem, correctly derived and implemented learners,
and an evaluation whose conclusions survive scrutiny.

\section{Dissertation structure and how to read it}
\label{sec:intro_structure}

The chapters follow the order in which the work must be understood, and
each builds strictly on its predecessor:

\begin{itemize}
  \item \textbf{Chapter~\ref{ch:background}} reviews the background: what
        an MDP is, how the two algorithm families work and where they
        came from, what is known about applying RL to trading, and the
        evaluation pitfalls this project is designed to avoid.
  \item \textbf{Chapter~\ref{ch:methodology}} is the mathematical core:
        the trading MDP is formulated precisely, and both learning rules
        are derived step by step from the objective --- including the
        complete answer to why a logarithm appears in the policy-gradient
        loss.
  \item \textbf{Chapter~\ref{ch:implementation}} shows how the
        mathematics becomes verified code: environment, agents, tests,
        and the exact hyperparameters used.
  \item \textbf{Chapter~\ref{ch:results}} reports the experiments in the
        order the methodology was built: game validation, the pilot
        collapse study, the main algorithm ablation, and the fee
        sensitivity study, with analysis of every finding.
  \item \textbf{Chapter~\ref{ch:conclusions}} closes with achievements
        against objectives, an honest critical appraisal, future work,
        and a reflection on project management.
  \item \textbf{Appendices} provide the equation-to-code traceability map
        (\ref{app:trace}), the complete per-month results
        (\ref{app:permonth}), and a plain-language glossary
        (\ref{app:glossary}).
\end{itemize}

A reader without a technical background can follow the storyline through
the chapter introductions, the plain-language summary that closes each
chapter, and the glossary; the mathematics in
Chapters~\ref{ch:methodology}--\ref{ch:implementation} deepens, but never
replaces, that narrative.

\paragraph{Chapter summary (plain language).} This project teaches two
classic learning algorithms, built from the ground up, to trade a stock
fund with imaginary money on real historical prices. Because markets
never tell you what the right decision was, the project's first rule is
honesty: prove the learners work on video games first, make them show
their working at every step, and compare them against the simplest
sensible strategy --- just buying and holding. The next chapter explains
the ideas this is built on.
